\bookchapter{第七章：前沿功能全解析}

\chapterintro{科技日新月异，手机功能也在不断进化。卫星通信、无网通、车祸检测、星闪技术……这些听起来很科幻的功能，其实已经走进我们的生活。让我们一起来了解这些前沿技术。}

\booksection{一、卫星通信：没信号也能求救}

卫星通信是近年来手机领域的一大突破。当你在荒郊野外、深山老林，手机完全没有信号时，卫星通信可能就是救命稻草。

\booksubsection{工作原理：}

手机直接与轨道上的通信卫星建立连接，发送求救信息。目前主要支持发送短信和位置信息，不能打电话。

\booksubsection{支持机型：}

\begin{itemize}
  \item \textbf{苹果}：iPhone 14系列及以后，支持通过卫星发送紧急SOS。iPhone 17系列进一步支持卫星发送iMessage和短信，且扩展了更多地区的覆盖。
  \item \textbf{华为}：Mate 50系列开始支持北斗卫星消息（仅文字），Mate 60 Pro系列进一步支持天通卫星通话（可打电话），Mate 70系列卫星通信能力进一步增强。
  \item \textbf{小米/OPPO/vivo}：2024年起旗舰机型陆续支持北斗卫星消息功能。
  \item \textbf{发展趋势}：卫星通信正从高端旗舰专属功能向下普及，未来中端机型也有望支持。中国北斗卫星系统在国产手机上的支持最为广泛。
\end{itemize}

\booksubsection{使用场景：}

\begin{itemize}
  \item 户外探险、登山遇险。
  \item 海上遇险。
  \item 自然灾害导致地面通信中断。
\end{itemize}

\booksubsection{使用方法：}

\begin{enumerate}
  \item 确保在开阔地带，能看到天空。
  \item 打开手机的卫星通信功能（通常在紧急呼叫界面）。
  \item 按照屏幕指引，对准卫星方向。
  \item 发送预设的求救信息和位置。
\end{enumerate}

\begin{tcbinfo}{小知识}
卫星通信需要手机有特殊的硬件天线，不是所有手机都支持。另外，卫星通信通常需要付费订阅服务，但紧急求救功能一般是免费的。

\end{tcbinfo}

\booksection{二、无网通：近距离通信新方式}

“无网通”是指在没有移动网络和Wi-Fi的情况下，两台手机之间直接通信的技术。

\booksubsection{技术原理：}

利用蓝牙或专用的无线电频率，在短距离内（通常几百米到几公里）建立点对点连接。

\booksubsection{应用场景：}

\begin{itemize}
  \item \textbf{户外活动}：徒步、露营时队友间通信。
  \item \textbf{灾难救援}：地震、洪水后通信中断时。
  \item \textbf{大型活动}：演唱会、体育赛事人多信号差时。
\end{itemize}

\booksubsection{支持平台：}

\begin{itemize}
  \item \textbf{苹果}：iOS 17的"NameDrop"和"离线找朋友"功能。
  \item \textbf{华为}：畅连App支持无网通话。
  \item \textbf{小米}：小米妙享中心的部分功能。
\end{itemize}

\booksection{三、车祸检测：关键时刻救命}

车祸检测是利用手机传感器检测严重车祸，并自动拨打紧急电话的功能。

\booksubsection{工作原理：}

手机通过加速度计、陀螺仪、气压计、麦克风等多个传感器，综合判断是否发生了严重碰撞。检测到车祸后，手机会：

\begin{enumerate}
  \item 弹出倒计时界面，询问你是否没事。
  \item 如果倒计时结束你没有响应，自动拨打紧急电话。
  \item 向紧急联系人发送你的位置。
\end{enumerate}

\booksubsection{支持机型：}

\begin{itemize}
  \item \textbf{苹果}：iPhone 14系列及以后。
  \item \textbf{安卓}：部分高端机型陆续支持。
\end{itemize}

\booksubsection{注意事项：}

\begin{itemize}
  \item 功能需要开启才能使用。
  \item 可能会误触发（比如过山车、剧烈运动），及时取消即可。
  \item 依赖传感器和算法，不是100\%可靠，但多一层保障。
\end{itemize}

\booksection{四、星闪技术：下一代无线连接}

星闪（SparkLink）是华为牵头研发的新一代无线通信技术，目标是在某些场景下替代蓝牙和Wi-Fi。它在关键指标上对标甚至超越现有方案：空口速率可达900Mbps（远超蓝牙的50Mbps），延迟低至20$\mu$s（蓝牙约15ms），单网络支持4096个设备同时连接（蓝牙通常8个以内）。

\booksubsection{技术特点：}

\begin{itemize}
  \item \textbf{超低延迟}：端到端延迟仅20微秒（0.02ms），比蓝牙低两个数量级，游戏和音乐同步几乎无感知。
  \item \textbf{高速率}：单链路峰值速率900Mbps，是蓝牙音频的数十倍，可传输无损音频甚至高清视频。
  \item \textbf{低功耗}：待机功耗与蓝牙相当，但传输效率更高，单位数据的能耗更低。
  \item \textbf{广覆盖+高并发}：连接距离比蓝牙远50\%以上，单网络支持上千设备同时在线。
\end{itemize}

\booksubsection{应用场景：}

\begin{itemize}
  \item \textbf{无线耳机}：更低延迟，玩游戏音画同步。
  \item \textbf{智能手表}：更稳定的连接。
  \item \textbf{智能家居}：更多设备同时连接。
  \item \textbf{车载互联}：更流畅的手机-车机连接。
\end{itemize}

\booksubsection{支持设备：}

目前主要是华为生态设备，如华为Mate 60系列、部分华为耳机和手表。未来会有更多品牌支持。

\textbf{技术展望}\par

星闪技术代表了无线通信的未来方向。随着更多厂商加入，我们可能会看到更多支持星闪的设备，享受更好的无线体验。

\begin{tcbpractice}{去检查一下：你的手机是否支持这些功能}
\begin{enumerate}
  \item 去设置里搜索"卫星"，看看是否有卫星通信相关选项。
  \item 搜索"车祸检测"或"紧急SOS"，检查是否开启。
  \item 如果你的手机是华为，看看是否支持星闪设备连接。
  \item 了解这些功能的使用方法，关键时刻可能派上用场。
\end{enumerate}

\end{tcbpractice}
